\chapter{Arquitectura del Software (Clean Architecture + CQRS + OWASP)}

\section{Patrón Arquitectónico: Clean Architecture en 4 Capas}
Para garantizar la mantenibilidad, escalabilidad y la independencia de frameworks o tecnologías externas, \textbf{SYSELL Web v2.0} adopta los principios de \textbf{Clean Architecture} propuestos por Robert C. Martin. La arquitectura se organiza en cuatro capas concéntricas con una regla de dependencia estricta: \textit{las dependencias del código apuntan únicamente hacia el interior}, asegurando que las reglas de negocio nunca dependan de bibliotecas, bases de datos o interfaces de usuario.

\begin{figure}[H]
\centering
\begin{tikzpicture}[
    scale=0.92, every node/.style={transform shape},
    cPres/.style={rectangle, draw=blue!80!black, fill=blue!10, rounded corners=6pt, line width=1.2pt, minimum width=12.8cm, minimum height=5.2cm},
    cInfra/.style={rectangle, draw=purple!80!black, fill=purple!10, rounded corners=5pt, line width=1.2pt, minimum width=10.6cm, minimum height=3.9cm},
    cApp/.style={rectangle, draw=cyan!80!black, fill=cyan!10, rounded corners=4pt, line width=1.2pt, minimum width=8.4cm, minimum height=2.7cm},
    cDom/.style={rectangle, draw=green!70!black, fill=green!15, rounded corners=3pt, line width=1.2pt, minimum width=6.2cm, minimum height=1.6cm, align=center}
]

% Capa 4: Presentacion
\node[cPres] (pres) at (0, 0) {};
\node[anchor=north, font=\bfseries\small, text=blue!80!black] at (0, 2.5) {Capa de Presentación (Controladores REST Laravel, React 18 SPA, Web Mobile)};

% Capa 3: Infraestructura
\node[cInfra] (infra) at (0, -0.3) {};
\node[anchor=north, font=\bfseries\small, text=purple!80!black] at (0, 1.6) {Capa de Infraestructura (MySQL 8.0, Sanctum Auth, Repositorios Eloquent, PDF/Excel)};

% Capa 2: Aplicacion
\node[cApp] (app) at (0, -0.6) {};
\node[anchor=north, font=\bfseries\small, text=cyan!80!black] at (0, 0.7) {Capa de Aplicación (Commands, Queries, Application Services, Interfaces)};

% Capa 1: Dominio
\node[cDom] (dom) at (0, -0.9) {\textbf{Capa de Dominio (Núcleo Puro)}\\\scriptsize Entidades (Prenda, Variante, Venta, Usuario), Value Objects (SKU, Moneda)\\\scriptsize Reglas Invariantes de Negocio (Sin dependencias externas)};

% Regla de Dependencia
\draw[-{Stealth[length=2.5mm]}, line width=1.8pt, red!80!black] (-5.8, -2.1) -- (-5.8, 2.1)
    node[midway, left=2pt, rotate=90, font=\scriptsize\bfseries, text=red!80!black] {Regla de Dependencia (Hacia el Dominio)};

\end{tikzpicture}
\caption{Diagrama de Cuatro Capas Concéntricas de Clean Architecture en SYSELL Web}
\label{fig:clean_architecture}
\end{figure}

\subsection{Descripción de Responsabilidades por Capa}
\begin{itemize}
    \item \textbf{Capa de Dominio:} Contiene la lógica esencial del negocio de ropa masculina. Las entidades como \texttt{Producto} o \texttt{Variante} hacen cumplir reglas invariantes (por ejemplo, el precio de venta debe ser estrictamente mayor a cero y el SKU debe cumplir con la nomenclatura \texttt{[CAT]-[MOD]-[COL]-[TAL]}). Esta capa no contiene dependencias de Laravel, React ni SQL.
    \item \textbf{Capa de Aplicación:} Implementa los casos de uso del sistema. Orquesta el flujo de información utilizando interfaces de repositorios abstractas (\texttt{IProductoRepository}, \texttt{IVentaRepository}), sin conocer si los datos se guardan en MySQL, PostgreSQL o memoria.
    \item \textbf{Capa de Infraestructura:} Implementa las interfaces de persistencia y servicios externos. Aquí residen las consultas SQL a MySQL 8.0, el almacenamiento de fotografías y el despacho de mensajes.
    \item \textbf{Capa de Presentación:} Provee los endpoints HTTP/JSON y renderiza los componentes interactivos de la interfaz web responsiva (Mobile-first).
\end{itemize}

\section{Patrón CQRS (Command Query Responsibility Segregation)}
Para optimizar el rendimiento y clarificar las intenciones transaccionales, se aplicó el patrón \textbf{CQRS}, separando conceptual y operativamente las acciones del sistema en dos modelos de ejecución:

\begin{figure}[H]
\centering
\begin{tikzpicture}[
    scale=0.88, every node/.style={transform shape},
    boxC/.style={rectangle, draw=red!70!black, fill=red!10, rounded corners=3pt, thick, align=center, minimum width=4.0cm, minimum height=1.1cm, font=\small},
    boxQ/.style={rectangle, draw=blue!70!black, fill=blue!10, rounded corners=3pt, thick, align=center, minimum width=4.0cm, minimum height=1.1cm, font=\small},
    dbNode/.style={cylinder, shape border rotate=90, draw=purple!80!black, fill=purple!15, aspect=0.3, minimum width=2.4cm, minimum height=2.4cm, align=center, font=\small\bfseries},
    flow/.style={thick, -{Stealth[length=2.2mm]}}
]

\node[rectangle, draw=gray!70, fill=gray!10, rounded corners=3pt, minimum width=2.8cm, minimum height=3.0cm, align=center] (client) at (-5.2, 0) {\textbf{Cliente Web / POS}\\\scriptsize (Peticiones HTTP)};

\node[boxC] (cmd) at (0, 1.4) {\textbf{Command Bus (Escritura)}\\\scriptsize Transacciones ACID (Venta, Ajuste)};
\node[boxQ] (qry) at (0, -1.4) {\textbf{Query Bus (Lectura)}\\\scriptsize Consultas Directas / Vistas};

\node[dbNode] (db) at (5.2, 0) {Base de Datos\\MySQL 8.0\\(3FN)};

% Flujos de Comandos
\draw[flow, red!80!black] (client.north) |- node[pos=0.72, above=2pt, font=\scriptsize\bfseries] {POST / PUT (Comandos)} (cmd.west);
\draw[flow, red!80!black] (cmd.east) -| node[pos=0.28, above=2pt, font=\scriptsize\bfseries] {Mutaciones ACID} (db.north);

% Flujos de Consultas
\draw[flow, blue!80!black] (client.south) |- node[pos=0.72, below=2pt, font=\scriptsize\bfseries] {GET (Catálogo, Reportes)} (qry.west);
\draw[flow, blue!80!black] (db.south) -| node[pos=0.28, below=2pt, font=\scriptsize\bfseries] {Proyecciones Optimizadas} (qry.east);

\end{tikzpicture}
\caption{Segregación Operativa de Comandos y Consultas con el Patrón CQRS}
\label{fig:cqrs_pattern}
\end{figure}

\begin{enumerate}
    \item \textbf{Commands (Comandos de Escritura):} Mutan el estado del sistema. Modifican datos, disparan transacciones ACID y retornan confirmación o IDs (ej. creación de prendas, confirmación de venta POS, anulación y ajustes de Kardex).
    \item \textbf{Queries (Consultas de Lectura):} Solo leen datos sin alterar el estado. Están optimizadas con proyecciones SQL directas y paginación (ej. consulta del catálogo web, stock por sucursal y reportería gerencial).
\end{enumerate}

\section{Seguridad Web (OWASP Top 10) y Autenticación Híbrida}
El sistema adopta las directrices de seguridad de \textbf{OWASP Top 10} para mitigar vulnerabilidades comunes:

\begin{table}[H]
\centering
\small
\begin{tabularx}{\textwidth}{>{\bfseries}p{3.8cm} >{\raggedright}p{4.8cm} Y}
\toprule
\textbf{Vulnerabilidad OWASP} & \textbf{Mecanismo de Protección en SYSELL} & \textbf{Implementación Técnica} \\
\midrule
A01: Broken Access Control & Control RBAC estricto a nivel de Middleware en cada endpoint API. & Roles \texttt{ADMINISTRADOR} y \texttt{VENDEDOR} evaluados en servidor antes de procesar la solicitud. \\
\addlinespace
A02: Cryptographic Failures & Cifrado unidireccional y canales seguros de transporte. & Contraseñas con hash Bcrypt (cost 12), forzando HTTPS con encabezados HSTS en producción. \\
\addlinespace
A03: Injection (SQLi) & Consultas parametrizadas obligatorias. & Uso exclusivo del Query Builder / Eloquent ORM; prohibido el concatenado de variables crudas en SQL. \\
\addlinespace
A07: Identification Failures & Bloqueo preventivo de cuentas y rate limiting. & Bloqueo automático al tercer fallo consecutivo (\texttt{RF-06}) y acelerador de retardo en reintentos. \\
\bottomrule
\end{tabularx}
\caption{Matriz de mitigación de vulnerabilidades de seguridad OWASP Top 10}
\label{tab:owasp_security}
\end{table}

\subsection{Estrategia Dual de Autenticación con Laravel Sanctum}
Para armonizar la seguridad del navegador con las necesidades de las terminales físicas de punto de venta:
\begin{enumerate}
    \item \textbf{Modo SPA Web (Navegador):} Utiliza cookies \texttt{HttpOnly}, \texttt{Secure} y \texttt{SameSite=Lax} con validación de token CSRF para el portal web administrativo y la tienda React, protegiendo las credenciales contra ataques de robo de sesión (XSS).
    \item \textbf{Modo API Token (Terminal POS / Sincronización):} Utiliza \textit{Bearer Tokens} (Personal Access Tokens) en el encabezado HTTP \texttt{Authorization} para permitir la operación en terminales de caja dedicadas y el envío de lotes de sincronización offline (\texttt{CU-10}).
\end{enumerate}
